\subsection{Group Relative Policy Optimization}
\label{sec:grpo}

The preceding sections traced a clear trajectory: PPO
(Section~\ref{sec:ppo}) demonstrated that reinforcement learning can
substantially improve LLM alignment but at the cost of training and serving
four models simultaneously---policy, reference, reward, and value networks.
DPO (Section~\ref{sec:dpo}) eliminated the reward model and value network
altogether but sacrificed on-policy exploration and the ability to leverage
verifiable scalar rewards.  \emph{Group Relative Policy Optimization}
(GRPO), introduced by \cite{shao2024deepseekmath} in the DeepSeekMath
project, charts an elegant middle course: it retains the on-policy sampling
and explicit reward signal of PPO while discarding the value network that
constitutes its principal computational bottleneck.  Because GRPO is the
algorithm at the heart of the scaling experiments presented in later
chapters, this section provides a detailed treatment of its derivation,
mechanics, and empirical motivation.

%% ---------- Motivation ----------
\paragraph{Motivation: the cost of the critic.}
Recall from Section~\ref{sec:ppo} that the PPO objective relies on an
advantage $A_t$ computed via Generalised Advantage Estimation
(GAE)~\cite{schulman2017ppo}, which in turn requires a \emph{learned value
function} $V_\psi$ of comparable size to the policy model $\pi_\theta$.
Training the value function alongside the policy introduces three
burdens:
\begin{enumerate}[leftmargin=1.8em,itemsep=2pt]
\item \textbf{Memory.}  The standard PPO setup for LLMs loads four
      full-sized models into GPU memory: the policy $\pi_\theta$, the old
      (frozen) reference policy $\pi_{\text{ref}}$, the reward model
      $r_\varphi$, and the value model $V_\psi$.  For a 7\,B-parameter
      model in mixed precision, each copy consumes roughly 14\,GB, so the
      aggregate footprint easily exceeds the capacity of a single
      high-end GPU.
\item \textbf{Training complexity.}  The value network must be trained
      concurrently with the policy, doubling the number of backward passes
      per update step and introducing an additional set of hyperparameters
      (value-function learning rate, GAE discount factor $\gamma$, GAE
      smoothing parameter $\lambda$).
\item \textbf{Credit-assignment difficulty.}  In the LLM setting, the
      reward model typically assigns a single scalar score to the
      \emph{entire} completion, yet GAE requires per-token value
      estimates.  Learning an accurate token-level value function from a
      single end-of-sequence reward is a challenging regression problem
      that can introduce high variance into the advantage estimates.
\end{enumerate}

\noindent
GRPO eliminates $V_\psi$ entirely.  Instead of learning a baseline for
variance reduction, it \emph{computes} a baseline on the fly by sampling
multiple completions for each prompt and normalizing their rewards
within the group.  This reduces the number of trained models from four to
three---policy, reference, and reward---cutting memory requirements by
approximately 25\% and simplifying the training loop.

%% ---------- Algorithm in detail ----------
\paragraph{The GRPO algorithm.}
GRPO proceeds in three conceptually distinct steps for each batch of
prompts.  We describe each step in turn and then present the full
pseudocode in Algorithm~\ref{alg:grpo}.

\subparagraph{Step 1: Group sampling.}
For each prompt $q$ drawn from the task distribution $\mathcal{D}$, the
current (old) policy generates a \emph{group} of $G$ independent
completions:
\begin{equation}
\label{eq:grpo-sampling}
\{o_1, o_2, \ldots, o_G\} \sim \pi_{\theta_{\text{old}}}(\cdot \mid q).
\end{equation}
Each $o_i = (o_{i,1}, o_{i,2}, \ldots, o_{i,|o_i|})$ is a variable-length
token sequence.  The group size $G$ is a key hyperparameter;
\cite{shao2024deepseekmath} use $G = 64$ in their experiments on
DeepSeekMath 7B.  Intuitively, a larger group provides a more reliable
estimate of the reward distribution for a given prompt, at the expense of
increased sampling cost.

\subparagraph{Step 2: Group-relative advantage estimation.}
A reward model $r_\varphi$ (or, more generally, any verifiable reward
function) scores each completion, yielding a vector of scalar rewards:
\begin{equation}
\label{eq:grpo-rewards}
\mathbf{r} = \{r_1, r_2, \ldots, r_G\},
\qquad r_i = r_\varphi(q, o_i).
\end{equation}
The key innovation of GRPO is that advantages are computed by
\emph{z-score normalization within the group}, without any learned value
function:
\begin{equation}
\label{eq:grpo-advantage}
\boxed{
\hat{A}_{i,t} = \tilde{r}_i
  = \frac{r_i - \operatorname{mean}(\mathbf{r})}
         {\operatorname{std}(\mathbf{r})},
\qquad i = 1, \ldots, G,
\quad t = 1, \ldots, |o_i|.}
\end{equation}
Under outcome supervision, every token in the $i$-th completion shares the
same advantage $\tilde{r}_i$; this is the setting adopted in the
DeepSeekMath experiments and in our own scaling study.  Note that the
advantage is \emph{relative}: a completion that happens to receive a
positive reward may still receive a negative advantage if most other
completions in the same group score higher.  This comparative structure
aligns naturally with reward models, which are themselves trained on
pairwise comparisons between outputs for the same prompt.

\subparagraph{Step 3: Clipped policy update with KL regularization.}
Given the group-relative advantages, GRPO updates the policy by
maximizing the following objective:
\begin{equation}
\label{eq:grpo-objective}
\begin{aligned}
\mathcal{J}_{\text{GRPO}}(\theta)
  &= \mathbb{E}_{\substack{q \sim \mathcal{D}\\
      \{o_i\}_{i=1}^{G} \sim \pi_{\theta_{\text{old}}}(\cdot|q)}}
    \left[
      \frac{1}{G}\sum_{i=1}^{G}
      \frac{1}{|o_i|}\sum_{t=1}^{|o_i|}
      \ell_{i,t}
    \right],\\
\ell_{i,t}
  &=
      \min\!\Bigl(
        \rho_{i,t}\,\hat{A}_{i,t},\;
        \operatorname{clip}\bigl(\rho_{i,t},\,
          1{-}\varepsilon,\,1{+}\varepsilon\bigr)\,\hat{A}_{i,t}
      \Bigr)
      - \beta\, D_{\mathrm{KL}}^{(i,t)} .
\end{aligned}
\end{equation}
where the per-token probability ratio is
\begin{equation}
\label{eq:grpo-ratio}
\rho_{i,t}
  = \frac{\pi_\theta(o_{i,t} \mid q,\, o_{i,<t})}
         {\pi_{\theta_{\text{old}}}(o_{i,t} \mid q,\, o_{i,<t})},
\end{equation}
and the per-token KL divergence from the reference policy is estimated via
the unbiased estimator of \cite{schulman2017ppo}:
\begin{equation}
\label{eq:grpo-kl}
D_{\mathrm{KL}}^{(i,t)}
  = \frac{\pi_{\text{ref}}(o_{i,t} \mid q,\, o_{i,<t})}
         {\pi_\theta(o_{i,t} \mid q,\, o_{i,<t})}
  - \log\frac{\pi_{\text{ref}}(o_{i,t} \mid q,\, o_{i,<t})}
             {\pi_\theta(o_{i,t} \mid q,\, o_{i,<t})}
  - 1.
\end{equation}
This estimator is guaranteed to be non-negative and avoids the need to
compute full distributional KL, which would be intractable over the
vocabulary.  The coefficient $\beta$ controls the strength of the
regularization; \cite{shao2024deepseekmath} set $\beta = 0.04$.  The
clipping parameter $\varepsilon$ (typically $0.2$) serves the same
stabilizing role as in PPO, preventing any single update from moving the
policy too far from $\pi_{\theta_{\text{old}}}$.

An important design choice distinguishes GRPO from PPO in its treatment of
the KL penalty.  In PPO-based RLHF, the KL penalty is typically folded
\emph{into the reward} at each token, i.e.,
$r_t = r_\varphi(q, o_{\leq t}) - \beta \log
\frac{\pi_\theta(o_t \mid q, o_{<t})}{\pi_{\text{ref}}(o_t \mid q, o_{<t})}$,
which then feeds into GAE to compute advantages.  GRPO instead adds the KL
term \emph{directly to the loss} (Eq.~\ref{eq:grpo-objective}), keeping the
advantage computation clean and free of reference-policy dependence.

%% ---------- Algorithm environment ----------
\begin{algorithm}[t]
\DontPrintSemicolon
\SetAlgoLined
\SetKwInOut{Input}{Input}
\SetKwInOut{Output}{Output}
\Input{Initial policy $\pi_{\theta_{\text{init}}}$;\;
       reward model $r_\varphi$;\;
       task prompts $\mathcal{D}$;\;
       hyperparameters $\varepsilon, \beta, G, \mu$}
\Output{Optimized policy $\pi_\theta$}
\BlankLine
$\pi_\theta \leftarrow \pi_{\theta_{\text{init}}}$ \;
\For{iteration $= 1, \ldots, I$}{
  $\pi_{\text{ref}} \leftarrow \pi_\theta$
    \tcp*{Sync reference model}
  \For{step $= 1, \ldots, M$}{
    Sample a mini-batch $\mathcal{D}_b \subset \mathcal{D}$ \;
    $\pi_{\theta_{\text{old}}} \leftarrow \pi_\theta$
      \tcp*{Snapshot policy}
    \ForEach{prompt $q \in \mathcal{D}_b$}{
      Sample $G$ outputs $\{o_i\}_{i=1}^{G} \sim
        \pi_{\theta_{\text{old}}}(\cdot \mid q)$
        \tcp*{Step 1}
      Compute rewards $\{r_i\}_{i=1}^{G}$ using $r_\varphi$ \;
      $\hat{A}_{i,t} \leftarrow
        \dfrac{r_i - \operatorname{mean}(\mathbf{r})}
              {\operatorname{std}(\mathbf{r})}$
        \quad $\forall\, i,t$
        \tcp*{Step 2}
    }
    \For{GRPO sub-iteration $= 1, \ldots, \mu$}{
      Update $\pi_\theta$ by maximizing $\mathcal{J}_{\text{GRPO}}(\theta)$
        \quad (Eq.~\ref{eq:grpo-objective})
        \tcp*{Step 3}
    }
  }
  \textbf{(Optional)} Update $r_\varphi$ via continual training with replay
    \tcp*{Iterative RL}
}
\Return $\pi_\theta$ \;
\caption{Group Relative Policy Optimization (GRPO)}
\label{alg:grpo}
\end{algorithm}

%% ---------- Comparison with PPO ----------
\paragraph{Comparison with PPO.}
Table~\ref{tab:ppo-vs-grpo} summarizes the principal differences between PPO
and GRPO in the context of LLM fine-tuning.  The most consequential
distinction is the replacement of the learned value function with
group-relative normalization.  In PPO, the advantage at token $t$ is
\begin{equation}
\label{eq:ppo-gae}
A_t^{\text{PPO}} = \sum_{l=0}^{T-t}(\gamma\lambda)^l\,
  \bigl(r_{t+l} + \gamma\, V_\psi(s_{t+l+1}) - V_\psi(s_{t+l})\bigr),
\end{equation}
which requires both the reward signal $r_{t+l}$ and per-state value
predictions from $V_\psi$.  GRPO replaces this entire computation with the
closed-form z-score of Eq.~\eqref{eq:grpo-advantage}, which depends only on
the rewards within the current group.  This substitution is particularly
well-suited to the LLM setting, where the reward model typically assigns a
single scalar to the full completion rather than per-token rewards, making
the training of an accurate token-level value function both difficult and
wasteful.

\begin{table}[t]
\centering
\caption{Comparison of PPO and GRPO for LLM reinforcement learning.}
\label{tab:ppo-vs-grpo}
\small
\begin{tabular}{@{}p{0.22\textwidth}p{0.36\textwidth}p{0.32\textwidth}@{}}
\toprule
\textbf{Aspect} & \textbf{PPO} & \textbf{GRPO} \\
\midrule
Models in memory &
  Policy, reference, reward, value (4) &
  Policy, reference, reward (3) \\
Advantage estimation &
  GAE via learned $V_\psi$ &
  Z-score within sampled group \\
KL regularization &
  Penalty added to reward &
  Penalty added directly to loss \\
Per-token value network &
  Required &
  Not needed \\
Sampling strategy &
  Single rollout per prompt &
  $G$ rollouts per prompt \\
Hyperparameters &
  $\varepsilon, \gamma, \lambda, \beta,$ value LR &
  $\varepsilon, \beta, G$ \\
\bottomrule
\end{tabular}
\end{table}

The elimination of the value network also yields a meaningful reduction in
wall-clock training time.  In the PPO loop, each update step requires a
forward pass through $V_\psi$ for every token in every rollout to compute
GAE, followed by a backward pass to update $V_\psi$---all in addition to
the policy gradient update.  GRPO replaces these operations with a single
reward-model inference per completion (not per token) and a trivial
mean/standard-deviation computation, resulting in faster iterations and
simpler engineering.

%% ---------- Reward design ----------
\paragraph{Reward design for mathematical reasoning.}
A distinctive feature of mathematical reasoning tasks is that correctness
is \emph{verifiable}: one can check whether a model's final answer matches
the ground truth without relying on a learned reward model.
\cite{shao2024deepseekmath} exploit this property by designing
\emph{rule-based} reward functions for their GRPO experiments:

\begin{itemize}[leftmargin=1.8em,itemsep=2pt]
\item \textbf{Accuracy reward.}  The primary signal is a binary indicator:
      $r_{\text{acc}}(q, o) = \mathbf{1}[\text{extract}(o) = a^*]$, where
      $a^*$ is the ground-truth answer and $\text{extract}(\cdot)$ is a
      deterministic parser that isolates the final boxed answer from the
      completion.  This reward is exact, noise-free, and incurs zero
      inference cost---a stark contrast to learned reward models, which
      are expensive to serve and susceptible to reward hacking.
\item \textbf{Format reward.}  An auxiliary reward encourages the model
      to produce well-structured chain-of-thought solutions.  For
      instance, a small bonus may be awarded for enclosing the final
      answer in a \verb|\boxed{}| environment or for following a
      prescribed step-by-step template.  Format rewards are typically
      combined additively with the accuracy reward:
      $r_i = r_{\text{acc}} + \alpha\, r_{\text{format}}$, where $\alpha$
      is a small weighting coefficient.
\end{itemize}

\noindent
The combination of verifiable binary rewards and group-relative
normalization creates a particularly clean learning signal.  Within a
group of $G = 64$ completions, some will be correct and others incorrect;
the z-score normalization assigns positive advantages to correct
completions and negative advantages to incorrect ones, with the magnitude
calibrated by the proportion of correct answers in the group.  When the
model is weak and most completions are wrong, the few correct ones
receive very large positive advantages, providing a strong gradient
signal.  As the model improves and more completions become correct, the
advantage magnitudes shrink, naturally annealing the learning rate.
This self-regulating property is an appealing feature of GRPO that
requires no manual scheduling.

%% ---------- Process supervision ----------
\paragraph{Extension to process supervision.}
While outcome supervision assigns a single reward per completion, GRPO
also supports \emph{process supervision}, in which a process reward model
scores each intermediate reasoning step.  Formally, if output $o_i$
consists of $K_i$ steps with end-of-step indices
$\text{idx}(1), \ldots, \text{idx}(K_i)$, the process reward model yields
step-level rewards $\{r^{\text{idx}(j)}_i\}_{j=1}^{K_i}$.  These are
again normalized via group-level statistics, and the advantage of each
token is set to the sum of normalized rewards from subsequent steps:
$\hat{A}_{i,t} = \sum_{\text{idx}(j) \geq t} \tilde{r}^{\text{idx}(j)}_i$.
Process supervision provides finer-grained credit assignment and can
improve sample efficiency on complex multi-step problems, though it
requires a more expensive process reward model.

%% ---------- Iterative RL ----------
\paragraph{Iterative GRPO.}
As training progresses, the policy may drift beyond the calibration range
of the initial reward model.  To address this, \cite{shao2024deepseekmath}
propose an \emph{iterative} variant (Algorithm~\ref{alg:grpo}, optional
step): after each outer iteration, the reward model is updated via
continual training on newly sampled completions from the current policy,
using a replay mechanism that retains 10\% of historical data to prevent
catastrophic forgetting.  The reference model is simultaneously reset to
the current policy checkpoint.  This co-evolution of reward model and
policy enables sustained improvement over multiple rounds of
reinforcement learning.

%% ---------- DeepSeekMath results ----------
\paragraph{DeepSeekMath results.}
Applied to the DeepSeekMath-Instruct 7B base model, GRPO training on
chain-of-thought instruction-tuning data from GSM8K and MATH yielded
striking gains.  On the competition-level MATH benchmark,
DeepSeekMath-RL 7B achieved a top-1 accuracy of \textbf{51.7\%}---a
4.9 percentage-point improvement over the supervised fine-tuned
baseline (46.8\%) and competitive with far larger proprietary models
such as Gemini-Ultra (53.2\%) and GPT-4 (52.9\%).  With
self-consistency decoding over 64 samples, accuracy rose further to
\textbf{60.9\%}.  On GSM8K, GRPO improved accuracy from 82.9\% to
88.2\%.  Crucially, these gains \emph{transferred to out-of-domain
benchmarks}: CMATH (Chinese mathematics) improved from 84.6\% to
88.8\%, even though no Chinese data were included in the RL training
set.  These results established GRPO as a practical and effective
alternative to PPO for mathematical reasoning, directly motivating
its adoption by the DeepSeek-R1 project~\cite{deepseek2025r1} for
training large-scale reasoning models.

\paragraph{Significance for this study.}
The appeal of GRPO for our scaling experiments is threefold.  First, the
elimination of the value network makes GRPO feasible on the GPU budgets
available to academic research groups---a consideration that is central to
the resource-constrained setting of this capstone project.  Second, the
algorithm's reliance on verifiable rewards rather than learned preference
models sidesteps the fragility and expense of reward-model training, an
important practical advantage when iterating rapidly across model sizes and
hyperparameter configurations.  Third, while
\cite{shao2024deepseekmath} demonstrated GRPO at a single model scale
(7B), \cite{deepseek2025r1} applied it to much larger models without
publishing detailed ablations of the scaling behaviour.  Our study aims to
fill this gap by systematically characterising how GRPO's training
dynamics, sample efficiency, and downstream accuracy evolve as a function
of model size, group size~$G$, KL coefficient~$\beta$, and
compute budget---questions that remain largely open in the literature.
